
\documentclass[11pt,a4paper]{article}
\usepackage[a4paper,margin=2.05cm]{geometry}
\usepackage[spanish,es-noquoting]{babel}
\usepackage[T1]{fontenc}
\usepackage[utf8]{inputenc}
\usepackage{lmodern,microtype}
\usepackage{amsmath,amssymb,amsthm,mathtools}
\usepackage{booktabs,longtable,array,enumitem}
\usepackage{xcolor,hyperref}
\pagecolor{white}\color{black}
\setlength{\parindent}{0pt}
\setlength{\parskip}{0.65em}
\newtheorem{teorema}{Teorema}
\newtheorem{observacion}[teorema]{Observación}
\title{\bfseries N57 --- Cierre Bockstein--Hensel del carry\\
\large \(P_m\) degradado, carry como ley, y blanco final APP/TPK}
\author{HMT / auditoría técnica}
\date{}
\begin{document}
\maketitle

\begin{abstract}
N57 cierra la pieza que había quedado amarilla en N56: los carries \(c_t,h_t\). El resultado es que no hay una palabra externa adicional que derivar para el margen; el carry es el cociclo Bockstein--Hensel inducido por la sección de dígitos \(\mathbb F_3\to\{0,1,2\}\subset\mathbb Z\). Se demuestra localmente en las \(27\) columnas posibles, se verifica la identidad de cociclo, se prueba que \(c,h\) son las obstrucciones no afines, y se confirma el caso terminal \(q=122120,\ c=000110\). El blanco final queda reducido a otro problema: derivar los estados de canal \(x_t^\chi,u_t^\chi\) desde APP/TPK.
\end{abstract}

\section{Cociclo local}

Para una columna:
\[
(p,e,f)\in\{0,1,2\}^3,
\]
definimos:
\[
q=(p+e+f)\bmod3,
\qquad
c=\frac{p+e+f-q}{3}.
\]
Para el canal firmado:
\[
\sigma=(1,1,-1),
\]
\[
a=(p+e-f)\bmod3,
\qquad
h=\frac{p+e-f-a}{3}.
\]
Con:
\[
\kappa(x,y)=\frac{x+y-((x+y)\bmod3)}{3},
\]
\[
\beta(x,y)=\frac{x-y-((x-y)\bmod3)}{3},
\]
y \(r=(p+e)\bmod3\), se obtiene:
\[
\boxed{
c=\kappa(p,e)+\kappa(r,f)
}
\]
\[
\boxed{
h=\kappa(p,e)+\beta(r,f).
}
\]
No hay parámetro libre.

\section{Identidad de cociclo}

La función \(\kappa\) satisface:
\[
\kappa(x,y)+\kappa((x+y)\bmod3,z)
=
\kappa(y,z)+\kappa(x,(y+z)\bmod3).
\]
La auditoría en las \(27\) ternas pasa:
\[
\boxed{
\delta\kappa=0.
}
\]
Esto es la forma local del Bockstein de la secuencia:
\[
0\to\mathbb Z\xrightarrow{\times3}\mathbb Z\to\mathbb F_3\to0.
\]

\section{Tabla universal}

\[
\begin{array}{ccc|rrrr}
p&e&f&q&c&a&h\\
\hline
0&0&0&0&0&0&0\\
0&0&1&1&0&2&-1\\
0&0&2&2&0&1&-1\\
0&1&0&1&0&1&0\\
0&1&1&2&0&0&0\\
0&1&2&0&1&2&-1\\
0&2&0&2&0&2&0\\
0&2&1&0&1&1&0\\
0&2&2&1&1&0&0\\
1&0&0&1&0&1&0\\
1&0&1&2&0&0&0\\
1&0&2&0&1&2&-1\\
1&1&0&2&0&2&0\\
1&1&1&0&1&1&0\\
1&1&2&1&1&0&0\\
1&2&0&0&1&0&1\\
1&2&1&1&1&2&0\\
1&2&2&2&1&1&0\\
2&0&0&2&0&2&0\\
2&0&1&0&1&1&0\\
2&0&2&1&1&0&0\\
2&1&0&0&1&0&1\\
2&1&1&1&1&2&0\\
2&1&2&2&1&1&0\\
2&2&0&1&1&1&1\\
2&2&1&2&1&0&1\\
2&2&2&0&2&2&0\\
\end{array}
\]

\section{No linealidad necesaria}

Tomando como diseño afín:
\[
1,\ p,\ e,\ f,
\]
se obtiene:
\[
q,a\text{ son afines sobre }\mathbb F_3.
\]
Pero:
\[
\boxed{
c,h\text{ no son afines.}
}
\]
Los tests de rango dan:
\[
\operatorname{rank}(A)=4,
\qquad
\operatorname{rank}(A|c)=5,
\qquad
\operatorname{rank}(A|h)=5.
\]
Por tanto:
\[
\boxed{
c,h=\text{obstrucción no lineal a la superposición de canales.}
}
\]

\section{Polinomios sobre \(\mathbb F_3\)}

La representación polinómica única sobre \(\mathbb F_3\) da:
\[
c\equiv 2ef + 2ef^2 + 2e^2f + 2pf + 2pf^2 + 2pe + pef + 2pe^2 + 2p^2f + 2p^2e\pmod3,
\]
\[
h\equiv 2f^2 + ef + 2ef^2 + e^2f + pf + 2pf^2 + 2pe + 2pef + 2pe^2 + p^2f + 2p^2e\pmod3.
\]
Esto muestra explícitamente que el carry es función algebraica interna, no dato externo.

\section{Ecuación Hensel global}

Para cada canal:
\[
y_t^\chi=x_t^\chi A,
\qquad
u_t^\chi=y_t^\chi\bmod3,
\qquad
x_{t+1}^\chi=\frac{y_t^\chi-u_t^\chi}{3}.
\]
Entonces:
\[
q_t=(u_t^\pi+u_t^e+u_t^\varphi)\bmod3,
\]
\[
c_t=\frac{u_t^\pi+u_t^e+u_t^\varphi-q_t}{3},
\]
\[
\boxed{
X_{t+1}^+
=
\frac{X_t^+A-q_t}{3}-c_t.
}
\]
Y para \(\alpha\):
\[
a_t=(u_t^\pi+u_t^e-u_t^\varphi)\bmod3,
\]
\[
h_t=\frac{u_t^\pi+u_t^e-u_t^\varphi-a_t}{3},
\]
\[
\boxed{
X_{t+1}^\sigma
=
\frac{X_t^\sigma A-a_t}{3}-h_t.
}
\]

\section{Terminal}

Para los bloques:
\[
021020,\quad001220,\quad100210,
\]
la auditoría local da:
\[
q=122120,
\qquad
c=000110.
\]
Además:
\[
r=221=2(1,1,-1).
\]
Esto coincide con el cierre terminal reportado.

\section{Dictamen}

\[
\boxed{
P_m\text{ no es primitivo.}
}
\]
\[
\boxed{
c_t,h_t\text{ quedan derivados como cociclos Bockstein--Hensel.}
}
\]
\[
\boxed{
\text{la superposición sin carry falla porque }c,h\text{ son obstrucciones no afines.}
}
\]
\[
\boxed{
\text{el blanco final ya no es carry: es APP/TPK}\to x_t^\chi,u_t^\chi.
}
\]
\end{document}
